\section{Implementation Details and Exploratory Experiments}
\label{app:implementation}

\subsection{Agent work directory}

Each \textsc{TVD}-Conflict run is executed in an isolated work directory that
holds the task specification, validator, active state, and retained conflicts.
A representative directory has the following structure:

\begin{verbatim}
run/<run-id>/
|-- TASK.md
|-- attempt.json
|-- judge.py
|-- validator.py
|-- successful.json
|-- agent.log
`-- logs/
    |-- <goal-id>_safe.jsonl
    `-- <goal-id>_unsafe.jsonl
\end{verbatim}

\texttt{TASK.md} defines the search objective, the safe and unsafe generation
strategies, the conflict criterion, and the candidate budget.
\texttt{attempt.json} is the exchange surface between the agent and validator:
it stores the current candidates and the feedback returned by the preceding
validation round. \texttt{judge.py} implements the shared binary and scoring
rubrics and sends identical visible content to both judges.
\texttt{validator.py} checks candidate structure, invokes the judges, evaluates
the conflict condition, and commits the resulting feedback.
\texttt{successful.json} contains the retained conflict pairs together with
both judgments. The files under \texttt{logs/} preserve per-goal trajectories
for the safe and unsafe branches, while \texttt{agent.log} records the agent
execution trace. This separation keeps candidate generation independent from
the validation rule and makes each retained pair traceable to its search
trajectory.

\subsection{Why the agent continues generating harmful jailbreak data}

The agent is not prompted as an end-user assistant answering a standalone
malicious request. It operates inside a structured research task in which its
role is to construct and revise evaluation artifacts, while a separate
validator determines whether each artifact satisfies the requested safety label
and induces judge disagreement. In the unsafe branch, harmful content is part
of the task-defined artifact rather than an incidental continuation; refusing
to construct it does not complete the assigned search. Validator feedback then
turns the task into a multi-step optimization process: the agent edits files,
submits candidates, observes two judgments, and revises unsuccessful cases.

This behavior is consistent with the Internal Safety Collapse mechanism
described by \citet{wu2026internal}: a capable model may continue producing
harmful content when such content becomes an internally required intermediate
or output of an otherwise structured task. In \textsc{TVD}-Conflict, the task
scaffold, persistent work state, and executable validation rule make that
requirement explicit. We use this connection as a mechanistic interpretation,
not as an independent causal test of Internal Safety Collapse.

\clearpage
\subsection{Cross-harness portability pilot}

We additionally ran an exploratory portability check using the Kimi CLI with
\texttt{kimi-code/k3-256k} at maximum reasoning effort. The pilot uses 50
AdvBench goals, safe and unsafe branches, a budget of ten rounds, and the
Qwen/DeepSeek judge pair. Each criterion therefore contains 100 target
instances. Table~\ref{tab:kimi-portability} reports the completed base-task
jobs recorded in \texttt{code/run/benchmark\_240\_me\_kimi/summary.csv}.

\begin{center}
  \refstepcounter{table}\label{tab:kimi-portability}
  \small
  \textbf{Table~\thetable:} Exploratory Kimi/K3 cross-harness results for the
  Qwen/DeepSeek judge pair. Goal coverage counts goals with at least one
  conflict in either branch; these results are excluded from the main
  experiment.

  \vspace{0.4em}
  \setlength{\tabcolsep}{5pt}
  \begin{tabular}{lrrrr}
    \toprule
    Criterion & Conflicts & CR@10 & Goal coverage & Mean rounds \\
    \midrule
    Binary & 92/100 & 92.0\% & 50/50 & 4.59 \\
    Scoring & 92/100 & 92.0\% & 50/50 & 4.14 \\
    \bottomrule
  \end{tabular}
\end{center}

At the recorded snapshot, 20 of the planned 80 shard-level jobs were complete;
the remaining configurations were still pending. The pilot therefore supports
cross-harness portability for this completed judge pair, but it is not a full
comparison between agent backbones and is kept separate from the 1,702 conflict
pairs reported in the main experiment.
